\chapter{Winkelverteilung und Pulshöhenspektrum der Höhenstrahlung}
\section{Theoretische Grundlagen}
\subsection{Höhenstrahlung}
\subsubsection{Primärkomponente}
\subsubsection{Sekundärkomponente}
\subsubsection{Winkelverteilung}
\subsection{Landau-Verteilung}
\section{Bestimmung der Schwellenspannung}
Um die Ereignisse über mehrere Tage messen zu können, muss zunächst die Diskriminatoren entsprechend eingestellt werden. Die Schwellenspannung wird so gewählt, dass Rauschen minimisiert wird, aber dennoch Muonen gemessen werden. 
\section{Durchführung der Messung}
\section{Auswertung}
\subsection{Winkelverteilung}
\subsection{Pulshöhenspektrum}